\chapter{Friends of friends}
\label{ch:friends}

\section{Triadic closure as a growth rule}

The previous chapter built networks by exploration. This one builds them by
introduction. Think about how you came to know the people you know. Only rarely
did you befriend a complete stranger; far more often a friend introduced you to
one of \emph{their} friends. This everyday mechanism --- a new tie forming
between two people who already share a common acquaintance --- has a name in
sociology: \emph{triadic closure}, because it ``closes'' a triangle by adding
the third side to a path of two. Two people who both know a third are, in a
sense, a triangle waiting to happen, and triadic closure is the event that
completes it. Sociologists have long seen this as the engine of social structure
\citep{granovetter1973}. If most new links close triangles, then triangles pile
up, and the network acquires the high cliquishness --- the large clustering
coefficient of the introduction --- that sets social networks apart from random
ones.

\begin{figure}[t]
\centering
\begin{adjustbox}{max width=\linewidth}\begin{tikzpicture}[
  nd/.style={circle,draw=black!65,fill=black!4,minimum size=6mm,inner sep=0pt},
  lnk/.style={draw=black!45,line width=0.8pt},
  pot/.style={draw=accent,line width=1.0pt,dashed},
  newl/.style={draw=accent,line width=1.5pt}]
  % panel (a)
  \begin{scope}
    \node[nd] (A) at (0,1.15)    {};
    \node[nd] (B) at (-1.2,-0.4) {};
    \node[nd] (C) at (1.2,-0.4)  {};
    \draw[lnk] (A)--(B);
    \draw[lnk] (A)--(C);
    \draw[pot] (B)--(C);
    \node[accent,font=\scriptsize] at (0,-0.72) {potential link};
    \node[font=\footnotesize] at (0,-1.5) {(a) friend of a friend};
  \end{scope}
  \node at (2.7,0.1) {\large $\Rightarrow$};
  % panel (b)
  \begin{scope}[xshift=4.9cm]
    \node[nd] (A2) at (0,1.15)    {};
    \node[nd] (B2) at (-1.2,-0.4) {};
    \node[nd] (C2) at (1.2,-0.4)  {};
    \draw[lnk] (A2)--(B2);
    \draw[lnk] (A2)--(C2);
    \draw[newl] (B2)--(C2);
    \node[font=\footnotesize] at (0,-1.5) {(b) triangle closed};
  \end{scope}
\end{tikzpicture}\end{adjustbox}
\caption{\textbf{Triadic closure.} Two people (bottom) who already share a common
friend (top) form a \emph{potential link} --- a friend-of-a-friend pair, drawn
dashed. Closing the triangle turns that potential link into a real one, the
elementary step by which the friends-of-friends rule builds social structure.}
\label{fig:fr-closure}
\end{figure}

Figure~\ref{fig:fr-closure} shows this elementary step. Notice that it is a
thoroughly \emph{local} rule. To introduce two friends you
need to know only one node and its neighbors: who you know, and who they know.
That is exactly the kind of local information this book is built on, and the
question, as always, is what such a rule produces in the aggregate. The answer,
worked out below following \citet{vazquez2003} and reviewed in
\citet{vazquez2025}, is that connecting friends-of-friends generates the whole
suite of social-network signatures at once: an effective ``rich get richer'' and
hence a power-law degree distribution, the clustering hierarchy in which
cliquishness falls with degree, and --- in pointed contrast to the technological
and biological networks of the other chapters --- \emph{assortative} degree
correlations, the tendency of well-connected people to know one another.

Models in this spirit have a history. \citet{davidsen2002} introduced an
acquaintance-network model (the DEB model, after its authors) in which people are
repeatedly introduced to friends of friends, producing a small world with high
clustering; \citet{holme2002} bolted an explicit triangle-forming step onto
preferential attachment to give it tunable clustering. What follows organizes
these ideas around a single bookkeeping device that makes the mechanism
transparent --- the \emph{potential link}.

\section{Potential links}

The key idea is that, under triadic closure, a pair of nodes is not just
``connected'' or ``not connected.'' A pair that already shares a common neighbor
is a triangle waiting to close --- a candidate in a way that an arbitrary pair of
strangers is not. So we sort every pair of nodes into one of three states:
\begin{itemize}
  \item \emph{disconnected} ($s$): no link and no common neighbor --- two
    strangers with no mutual acquaintance;
  \item connected by a \emph{potential link} ($p$): no link yet, but at least one
    common neighbor --- two friends-of-a-friend, poised to be introduced;
  \item connected by a \emph{link} ($e$): already friends.
\end{itemize}
The \emph{potential link} is the formal stand-in for ``friend of a friend.'' We
write $k_i$ for the ordinary degree of node $i$ (its number of friends) and
$k^*_i$ for its \emph{potential degree}, the number of friends-of-friends it has
--- the number of potential links waiting at it.

Triadic closure is then a story about a pair changing state: a potential link
($p$) becomes a real link ($e$) when the triangle closes. To track this we need
to know the rates at which pairs hop between the three states, and one
relationship among those rates is the mechanical heart of the model. It says
that potential links are born in proportion to a node's degree --- introduce
yourself to someone with many friends and you instantly acquire many
friends-of-friends. The next box sets up the bookkeeping and draws out its main
consequence; the conclusion is restated plainly afterward, so the box may be
skipped.

\begin{calculation}{from the three-state bookkeeping to ``rich get richer''}
Let $\nu_{x\to y}$ be the rate (per unit of network size $n$) at which a pair
hops from state $x$ to state $y$. Ignoring fluctuations, the average degree and
potential degree of node $i$ evolve as
\begin{align}
  n\,\frac{\partial k_i}{\partial n}
    &= \nu_{s\to e}\,\hat k_i + \nu_{p\to e}\,k^*_i
       - (\nu_{e\to s}+\nu_{e\to p})\,k_i, \nonumber\\
  n\,\frac{\partial k^*_i}{\partial n}
    &= \nu_{s\to p}\,\hat k_i + \nu_{e\to p}\,k_i
       - (\nu_{p\to s}+\nu_{p\to e})\,k^*_i,
  \label{eq:fr-master}
\end{align}
where $\hat k_i = n - k_i - k^*_i$ is the number of nodes that are neither
friends nor friends-of-friends of $i$. The link between the two kinds of
connection comes from the geometry of triangles: when a new node $i$ attaches to
an existing node $j$, it instantly shares the common neighbor $j$ with \emph{all}
of $j$'s neighbors, gaining a potential link to each. Hence potential links are
spawned in proportion to degree,
\begin{equation}
  \nu_{s\to p} = \nu_{s\to e}\, k_i, \qquad
  \nu_{p\to s} = \nu_{e\to s}\, k_i.
  \label{eq:fr-coupling}
\end{equation}
Now make two modeling choices. First, neglect link deletion, $\nu_{e\to s}=0$ ---
accurate for, say, a co-authorship network, where having written a paper together
is a fact that cannot later become false. Second, make closing a triangle far
likelier than linking two strangers,
\begin{equation}
  \nu_{s\to e} = \frac{\mu_0}{n^2}, \qquad
  \nu_{p\to e} = \frac{\mu_1}{n},
  \label{eq:fr-rates}
\end{equation}
the stranger-linking rate carrying an extra factor of $1/n$. Substituting these
into Eq.~\eqref{eq:fr-master} collapses the system to
\begin{equation}
  n\,\frac{\partial k_i}{\partial n} = \mu_0 + \mu_1\, k^*_i, \qquad
  n\,\frac{\partial k^*_i}{\partial n} = \mu_0\, k_i - \mu_1\, k^*_i.
  \label{eq:fr-reduced}
\end{equation}
\end{calculation}

The upshot of Eq.~\eqref{eq:fr-reduced} is that a node's degree grows at a rate
proportional to how many friends-of-friends it has --- and a node with many
friends has, through them, many friends-of-friends. So once again the rich get
richer, but this time the preference is manufactured by triadic closure rather
than by surfing: a well-connected person is simply involved in more potential
introductions. As in the previous chapter, no node ever consulted a global
ranking; the bias toward hubs is a free by-product of a purely local act.

\section{Degree distribution and clustering hierarchy}

Because the growth rate is linear in degree, the same reasoning that gave a power
law in the last chapter applies here, and it simultaneously fixes the
cliquishness. The box carries out both; the two results and their meaning are in
the paragraph after it.

\begin{calculation}{the power law and the clustering hierarchy}
Integrating the linear system~\eqref{eq:fr-reduced} gives power-law growth of
both degrees with the age of the node,
\begin{equation}
  k_i(n) = k_0 \left(\frac{n}{n_i}\right)^{\beta}, \qquad
  k^*_i(n) = k^*_0 \left(\frac{n}{n_i}\right)^{\beta},
  \label{eq:fr-growth}
\end{equation}
where $n_i$ is the network size when node $i$ arrived and
\begin{equation}
  \beta = \frac{\mu_1}{2}\left(-1 + \sqrt{1 + 4\,\frac{\mu_0}{\mu_1}}\,\right).
  \label{eq:fr-beta}
\end{equation}
With nodes added at a constant rate, $P(n_i{=}n)=1/n$, a node has degree above $k$
only when it is old enough for Eq.~\eqref{eq:fr-growth} to exceed $k$; averaging
over its birth time gives the cumulative distribution
\begin{equation}
\adjustbox{max width=\linewidth}{$\displaystyle
  P(k_i > k) = P\!\left[k_0\left(\frac{n}{n_i}\right)^{\beta} > k\right]
  = \int_0^n \frac{dn_i}{n}\,
    \Theta\!\left[k_0\left(\frac{n}{n_i}\right)^{\beta} - k\right]
$}
  \label{eq:fr-cumulative}
\end{equation}
with $\Theta[\cdot]$ the step function, $1$ when its bracket holds and $0$
otherwise. Differentiating it leaves a power-law degree distribution,
\begin{equation}
  p_k = \frac{\partial P(k_i > k)}{\partial k} \sim k^{-\gamma}, \qquad
  \gamma = 1 + \frac{1}{\beta}.
  \label{eq:fr-pk}
\end{equation}
For the clustering, the dominant way node $i$ gains a triangle is a potential
link closing: when the potential link from $i$ to $j$ (sharing neighbor $k$)
becomes real, $i$ gains both the neighbor $j$ and a link $j$--$k$ among its
neighbors. Thus $\partial e_i/\partial n = \nu_{p\to e}\,k^*_i = \mu_1 k^*_i/n$,
and integrating with Eq.~\eqref{eq:fr-growth},
\begin{equation}
  \langle c\rangle_k = \frac{2 e(k)}{k(k-1)} \approx \frac{2\mu_1}{k}.
  \label{eq:fr-ck}
\end{equation}
\end{calculation}

Two familiar results emerge. First, Eq.~\eqref{eq:fr-pk}: the degree
distribution is again a power law. It is worth knowing exactly where the power
law comes from --- it comes from the proportionality in
Eq.~\eqref{eq:fr-coupling}, the fact that potential links are spawned in
proportion to degree. Were they instead spawned at a flat rate, the mechanism
would still create triangles but the degree distribution would decay
exponentially, with no fat tail of hubs. Scale-freeness is the fingerprint of
that degree-proportional spawning. Second, Eq.~\eqref{eq:fr-ck}: cliquishness
again falls as one over the degree --- the clustering hierarchy of the
introduction, now arriving by a third route (after search, and before
duplication in the next chapter). Here its origin is especially vivid: every
triangle a node owns is just a friend-of-a-friend pair that has been introduced.

\section{Connecting nearest neighbors}

The continuum theory is borne out by a minimal simulation, a growing version of
the DEB rule \citep{davidsen2002}. Starting from a single node and no links,
repeat:
\begin{itemize}
  \item with probability $1-u$, \emph{add} a new node and link it to a randomly
    chosen node $j$ --- which by definition creates a potential link between the
    newcomer and every neighbor of $j$;
  \item with probability $u$, \emph{close} a triangle: pick a potential link at
    random and turn it into a real link.
\end{itemize}
The single dial $u$ is the rate of triangle-closing relative to adding newcomers.
Feeding the mapping $\mu_0=1$, $\mu_1=u/(1-u)$ into the formulas above gives the
exponent in closed form, as the short box records.

\begin{calculation}{the exponent as a function of the closure rate}
With $\mu_0=1$ and $\mu_1=u/(1-u)$, Eqs.~\eqref{eq:fr-beta} and~\eqref{eq:fr-pk}
give
\begin{equation}
  \gamma(u) = 1 + \frac{2(1-u)}{u}
    \left(-1 + \sqrt{1 + 4\,\frac{1-u}{u}}\,\right)^{-1},
  \label{eq:fr-gammau}
\end{equation}
with limiting values
\begin{equation}
  \gamma(0) = \infty, \qquad \gamma(1) = 2.
  \label{eq:fr-limits}
\end{equation}
\end{calculation}

In words, the exponent is tunable by how eagerly the network closes triangles.
When closure is frequent and newcomers rare ($u\to 1$) the exponent sinks to its
minimum of $2$, the most hub-dominated case; when newcomers arrive so fast that
triangles never get a chance to close ($u\to 0$) the exponent blows up and the
heavy tail of hubs is lost. Simulations up to a million nodes confirm the
power-law degree distribution (the mean-field formula slightly overestimating the
exponent, as in the previous chapter), and they confirm the clustering hierarchy,
with cliquishness falling as a power of degree.

\section{Why social networks are assortative}

The most telling prediction is about the degree correlations --- and here the
friends-of-friends rule parts company with every other rule in this book.
Measuring the average neighbor degree $\langle k_{nn}\rangle(k)$ in the simulated
network, one finds it \emph{rises} with degree: the well-connected are connected
to the well-connected. The network is \emph{assortative}
\citep{newman2002} --- ``birds of a feather'' --- exactly as measured in real
collaboration and co-authorship networks, and exactly opposite to the
disassortative Internet and protein networks of the introduction.

Why the opposite sign? The reason becomes obvious once seen. Every new link here
closes a triangle, joining two nodes that share a common neighbor. But a node
that serves as the common neighbor of many pairs is, by definition, a high-degree
node, and the two nodes it brings together are drawn from its own neighborhood,
which is better connected than the network at large. So closure tends to wire
together nodes of \emph{similar}, and similarly high, degree. The very mechanism
that the disassortative models of the other chapters lack --- linking within a
neighborhood rather than reaching out to a random or peripheral node --- is what
tips the correlations positive. This is the deeper lesson of the chapter, and of
the review that develops it \citep{vazquez2025}: the \emph{sign} of the degree
correlations is not an arbitrary fitting knob but a direct readout of the
microscopic rule. Search and duplication reach outward, to the periphery, and so
make hubs that shun one another; triadic closure works inward, within
neighborhoods, and so makes hubs that keep each other's company.

\section{Summary}

Modeling social growth by triadic closure --- a newcomer attaches to a node and
then, over time, to that node's friends --- and formalizing ``friend of a friend''
as a \emph{potential link}, we obtain a tractable theory in which potential links
are created in proportion to degree. From that one ingredient follow an effective
``rich get richer,'' a power-law degree distribution with a tunable exponent
$\gamma = 1 + 1/\beta$, a clustering hierarchy in which cliquishness falls as one
over the degree, and --- distinctively --- assortative degree correlations.
Connecting friends-of-friends thus reproduces, from purely local moves, the
characteristic shape of social networks, and explains why the sign of their
correlations differs from that of the technological and biological networks
treated elsewhere in this book.
